% TOP_PROBLEM: {"id": 20000958, "title": "A cut-cone certificate for binary transitive hosts", "category_id": 2, "category": {"id": 2, "name": "combinatorics", "display_name": "Combinatorics", "description": "Counting problems, graph theory, discrete structures.", "slug": "combinatorics", "order_index": 2, "created_at": "2026-07-31T15:26:25.671Z"}, "status": "partially_solved", "problem_number": "AIM-COMBINATORICS-0083", "set_id": 14}
% FIRST_POSTED: 2026-10-01
% ENTRY_KIND: statement_only
% UnsolvedMath ID 20000958; source catalogue code AIM-COMBINATORICS-0083
% !TeX program = lualatex
% Definition and sources only; this file is not an LLM solution attempt.
\documentclass[11pt,a4paper]{article}
\usepackage[margin=25mm]{geometry}
\usepackage{fontspec}
\setmainfont{lmroman10-regular.otf}[BoldFont=lmroman10-bold.otf,
 ItalicFont=lmroman10-italic.otf,BoldItalicFont=lmroman10-bolditalic.otf]
\usepackage{amsmath,amssymb,mathtools}
\usepackage{unicode-math}
\setmathfont{latinmodern-math.otf}
\AtBeginDocument{\let\setminus\smallsetminus}
\usepackage{xurl,hyperref}
\hypersetup{colorlinks=true,urlcolor=blue}
\setcounter{secnumdepth}{0}
\setlength{\parindent}{0pt}
\setlength{\parskip}{5pt}
\setlength{\emergencystretch}{3em}
\begin{document}
\section*{A cut-cone certificate for binary transitive hosts}\label{20000958}
{\small UnsolvedMath ID 20000958 · AIM-COMBINATORICS-0083 · Combinatorics · AIM Workshop Problem Lists}
\subsection{Definitions and mathematical statement}
Let $X$ be a nonempty finite set of points in a Euclidean space
$\mathbb R^d$, with distances measured by the usual norm. Call $X$
Ramsey if for every integer $r\ge1$ there is an integer $N\ge d$ such
that every coloring $c:\mathbb R^N\to\{1,\ldots,r\}$ admits an
isometric embedding $f:X\to\mathbb R^N$ with $c$ constant on $f(X)$.
An isometric embedding preserves every pairwise distance. No
measurability or regularity is imposed on the coloring, and the copy
must have exactly the original scale of $X$.

Call a finite set $Y\subseteq\mathbb R^m$ transitive if its Euclidean
symmetry group acts transitively on its points: for every $y,z\in Y$
there is an isometry $g$ of $\mathbb R^m$ satisfying $g(Y)=Y$ and
$g(y)=z$. Say $X$ is subtransitive if some such finite transitive $Y$
contains an isometric copy of $X$. The ambient dimension $m$ of this
host is allowed to differ from $d$.

The Leader--Russell--Walters conjecture asks whether a nonempty finite
Euclidean configuration is Ramsey if and only if it is subtransitive.
Thus both directions are part of the characterization. In the host
condition $Y$ must be finite, and its transitivity concerns the symmetry
group preserving $Y$, rather than arbitrary motions of the ambient
Euclidean space. The host may depend on $X$ but not on a later choice
of coloring.
\subsection{Short English statement}
Are the finite Euclidean configurations that force monochromatic congruent copies exactly those contained in finite transitive configurations?
\subsection{Sources}
\begin{itemize}
\item[S1] ULAM.AI and the UnsolvedMath Contributors, supplied problems.json, record 20000958 (AIM-COMBINATORICS-0083), AIM Workshop Problem Lists. Catalogue identity and source transcription; CC BY 4.0.
\url{https://huggingface.co/datasets/ulamai/UnsolvedMath}.
\item[S2] American Institute of Mathematics, Graph Ramsey theory, item 36. Original workshop question underlying AIM-COMBINATORICS-0083.
\url{https://www.overleaf.com/read/mnvcscjjysvg}.
\end{itemize}
\end{document}
